% Preamble patches for a MinerU OCR project translated to Chinese
% =============================================================
%
% MinerU's stock preamble already loads fontspec, polyglossia and ctex, and
% picks a CJK font through an \IfFontExistsTF chain.  The edits below are the
% ones that are easy to get wrong because *where* they go matters.
%
% Copy the relevant lines rather than adopting this file wholesale — several of
% these only make sense in the context of MinerU's existing preamble.


% --- 1. CJK fonts ---------------------------------------------------------
% Put the platform's own font first in MinerU's chain so the document builds on
% a machine without Source Han / Noto CJK installed.  On macOS "Songti SC" and
% "Heiti SC" are present out of the box and look right for a technical report.
%
% \setCJKmainfont and \setCJKsansfont are *preamble-only*: placing them after
% \begin{document} raises "Can be used only in preamble."  (The merged file
% begins the body with the centred title block, so it is easy to slide these in
% one line too late.)
\IfFontExistsTF{Songti SC}
{\newfontfamily\chinesefont{Songti SC}}
{\IfFontExistsTF{Source Han Serif CN}
{\newfontfamily\chinesefont{Source Han Serif CN}}
{\IfFontExistsTF{Noto Serif CJK SC}
  {\newfontfamily\chinesefont{Noto Serif CJK SC}}
  {\IfFontExistsTF{SimSun}
    {\newfontfamily\chinesefont{SimSun}}
    {\newfontfamily\chinesefont{Arial Unicode MS}}
}}}}
\setCJKmainfont{Songti SC}
\setCJKsansfont{Heiti SC}


% --- 2. URL line breaking -------------------------------------------------
% OCR often turns a wrapped URL into a URL with a space inside it.  Joining that
% space back is correct, but it can leave a 150-180pt unbreakable token that
% then shoots past the right margin.  Load xurl *before* hyperref so the URL
% can break at any character.
\usepackage{xurl}
\usepackage{hyperref}


% --- 3. Looser line breaking ---------------------------------------------
% With long URLs and long technical identifiers, a little extra stretch removes
% most residual overfull boxes.  Both belong in the preamble.
\sloppy
\setlength{\emergencystretch}{3em}


% --- 4. Wrappable table columns ------------------------------------------
% MinerU emits every table with `l` columns, and `l` never wraps.  Define a
% ragged-right p-column so long benchmark names and model names can wrap.
% Requires the array package (MinerU already loads it) and ragged2e.
\usepackage{ragged2e}
\newcolumntype{L}[1]{>{\RaggedRight\arraybackslash}p{#1}}


% --- 5. Centred title block ----------------------------------------------
% MinerU emits the paper title as \section{...}.  Replace it with a centred
% block.  This one is *document body* content, not preamble, but it is grouped
% here because it is part of the same frontmatter edit.
%
% \begin{center}
% {\LARGE DeepSeek-V4.1-Flash\par}
% \vspace{0.5em}
% {\large 突破 KV 缓存压缩的极限\par}
% \vspace{0.8em}
% {\normalsize DeepSeek-AI\quad research@deepseek.com\par}
% \end{center}


% --- 6. Caption names and the table of contents ---------------------------
% ctex does not relabel \caption prefixes; do it explicitly.  \figurename and
% \tablename may sit in the preamble.
\renewcommand{\figurename}{图}
\renewcommand{\tablename}{表}

% \contentsname is different: polyglossia re-activates the language after the
% preamble and resets it, so a preamble-level \renewcommand is silently lost.
% It must go in the body, immediately before \tableofcontents.
%
% \renewcommand{\contentsname}{目录}
% \tableofcontents
%
% When merging chunks, insert the TOC *after* the centred title block, not
% directly after \begin{document}: the preamble chunk in this pipeline carries
% \begin{document} and the title block follows it, so inserting at
% \begin{document} puts the contents pages ahead of the title page.
